\begin{table}[H]
\small
\setlength{\tabcolsep}{5pt}
\renewcommand{\arraystretch}{0.95}

\caption{\label{tab:country-heterogeneity-cumulative-gap}Cumulative inflation gap between the first and fifth income quintiles, June 2021--June 2023}
\centering
\begin{tabular}[t]{lllll}
\toprule
Country & Cumulative inflation & Q1 cumulative inflation & Q5 cumulative inflation & Inflation gap\\
\midrule
Austria & 17.2 & 17.5 & 17.1 & 0.4\\
Cyprus & 12.0 & 12.7 & 12.0 & 0.7\\
Estonia & 33.0 & 38.7 & 31.2 & 7.5\\
Finland & 12.5 & 11.5 & 12.7 & -1.2\\
France & 12.2 & 12.3 & 12.1 & 0.2\\
\addlinespace
Germany & 15.6 & 16.5 & 15.2 & 1.3\\
Greece & 14.7 & 15.1 & 14.5 & 0.6\\
Ireland & 14.9 & 17.2 & 13.3 & 3.8\\
Italy & 15.8 & 17.2 & 15.1 & 2.1\\
Latvia & 28.9 & 36.0 & 26.5 & 9.5\\
\addlinespace
Lithuania & 30.4 & 33.6 & 29.0 & 4.6\\
Luxembourg & 11.4 & 11.0 & 11.1 & -0.2\\
Netherlands & 16.9 & 17.7 & 16.5 & 1.2\\
Portugal & 14.2 & 13.1 & 15.4 & -2.3\\
Slovakia & 25.4 & 26.7 & 24.0 & 2.7\\
\addlinespace
Spain & 11.8 & 11.8 & 11.7 & 0.1\\
\midrule
\textbf{Euro area} & \textbf{14.7} & \textbf{15.2} & \textbf{14.4} & \textbf{0.9}\\
\bottomrule
\end{tabular}
\vspace{-0.4em}
\begin{flushleft}
\footnotesize{Notes. Entries are points and are rounded to one decimal; gaps are calculated from unrounded indices and may therefore differ slightly from the difference between the displayed group values. For every household category, cumulative inflation is calculated as $100(P_{g,2023\text{-}06}/P_{g,2021\text{-}06}-1)$, and the gap is the first group's cumulative inflation minus the second group's. The euro-area Q1--Q5 gap is 15.23 minus 14.36, or 0.86 point, rounded to 0.9. Applying the same endpoints, chaining, 2020 index base, and annual HICP country weights gives an age gap of 15.17 minus 13.90, or 1.27 points, rounded to 1.3, for households aged 60 or over minus those under 30; and a residence gap of 15.27 minus 14.31, or 0.96 point, rounded to 1.0, for rural areas minus cities. Q1 and Q5 cumulative inflation are computed from income-quintile price indices using RAS expenditure weights with the updated HBS--HICP bridge for imputed rents. France is computed at COICOP level 2, consistently with the other country rows; Germany and Portugal use successive Eurostat HBS waves, while Slovakia uses a homogeneous 2015 wave. The euro-area row is rebuilt from country-level indices and EA20 HICP country weights across all 20 member countries. Sources: Eurostat, authors' calculations.}
\end{flushleft}
\end{table}
